\scanpage{140}
$S$ schéma normal intègre, pt générique $\eta$, $k(\eta)=K$.

\uline{Isomotif} ($\frac12$ simple) \uline{lisse} sur $S$ : = isomotif semi-simple \inserted{$M$} sur $S$, tel que pour tout $\ell\in\mathbb P$, $\ell\ne\mathrm{car}\,K$, $T_\ell(M)$ soit non ramifié rel. à $S_\ell=S[1/\ell]$

Conjecturalement, il suffit que ce soit vrai pour \uline{un} $\ell$ premier à toutes les car. résiduelles ($\Rightarrow$ il suffit que ce soit vrai pour \uline{deux} $\ell\ne\ell'$).

Motif (\unsure{isosemisimple|iso semisimple}) lisse sur $S$ (si $\mathrm{car}\,K=0$) :
\[
 =\{M,\{\mathbf{E}_\ell\},\{\alpha_\ell\}\}
\]
où

a) $M$ un isomotif lisse sur $S$

b) $\forall\,\ell$, $E_\ell$ un faisceau $\ell$-adique et tordu sur $S_\ell=S[1/\ell]$ (ou encore sur $K$, mais non ramifié rel à $S[1/\ell]$)

c) $\forall\,\ell$, $\alpha_\ell:E_\ell\deleted{\otimes_{\mathbb Z_\ell}\mathbb Q_\ell}\xrightarrow{\sim}T_\ell(M)$

avec un axiome : ``pour presque tout $\ell$, $u_\ell$ provient d'un isomorphisme sur $\mathbb Z_\ell$ -- pas seulement à torsion près''.

Ceci implique que pour \unsure{presque|\deleted{presque}} tout $\ell$, $E_\ell$ est sans torsion.

Homomorphismes de motifs lisses sur $S$
\[
 (M,\{E_\ell\}\deleted{\unsure{u\to M}},\{\alpha_\ell\})\ \unsure{\mathord{-}|\longrightarrow}\ (M',\{E'_\ell\},\{\alpha'_\ell\})
\]
donné par $(u,\{u_\ell\})$, avec $u:M-M'$ un hom de isomotifs\deleted{\unsure{de}}, et $u_\ell:E_\ell-E'_\ell$ un hom de faisceaux $\ell$-adiques tels que pour tt $\ell$, on ait commutativité de
\diagram{\begin{tikzcd}[ampersand replacement=\&]
 E_\ell\otimes_{\mathbb Z_\ell}\mathbb Q_\ell \arrow[r,dash,"\alpha_\ell"] \arrow[d,dash,"u_\ell\otimes_{\mathbb Z_\ell}\mathbb Q_\ell"] \& T_\ell(M)\arrow[d,"T_\ell(u)"]\\
 E'_\ell\otimes_{\mathbb Z_\ell}\mathbb Q_\ell \arrow[r,dash,"\alpha'_\ell"'] \& T_\ell(M')\text{ .}
\end{tikzcd}}

\uline{N.B.} Le foncteur \deleted{\unsure{Hom}} \inserted{Motifs} sur $K$ $\longleftarrow$ Motifs lisses sur $S$ est pleinement fidèle. Les Hom sont des modules de type fini sur $\mathbb Z$. La catégorie des motifs lisses sur $S$ est abélienne (on prend \uline{noyaux}, conoyaux séparément pour les composantes $M$, $E_\ell$). L\deleted{e foncteur}\inserted{'image des} \deleted{Mot} (Motifs lisses sur $S$) \unsure{en|dans|$\to$} (Motifs sur $K$) est une ss-catégorie abélienne \uline{épaisse}.

Le foncteur $(\text{Motifs}\inserted{\text{ lisses}})\to(\text{Isomotifs}\inserted{\text{ lisses}})$ est exact ess\textsuperscript{t} surjectif
% archive sheet number: 139
